\chapter{Infrastruttura di calcolo}
\label{cap:infrastruttura}

Questo capitolo documenta l'ambiente di esecuzione e, soprattutto, cinque guasti
reali che ne sono derivati. I guasti sono riportati per esteso perché sono il
prodotto più trasferibile del progetto: nessuno di essi era un errore di
algoritmo, tutti erano errori di \emph{ordine} o di \emph{ambiente}, e tutti
producevano fallimenti silenziosi o diagnostiche fuorvianti.

\section{La catena di esecuzione}

L'addestramento non gira sulla macchina di sviluppo, ma su un cluster gestito da
SLURM. Il percorso completo è:

\begin{equation}
  \texttt{login} \;\to\; \texttt{srun} \;\to\; \texttt{compute}
  \;\to\; \texttt{apptainer} \;\to\; \texttt{setup}
\end{equation}

\begin{description}
  \item[login] Il nodo di accesso. Non ha il runtime del progetto: nessuna GPU,
        nessun ambiente Python con le dipendenze. Qualunque comando che importi
        il codice del progetto fallisce qui.
  \item[srun] Il comando di allocazione SLURM, che ottiene un nodo di calcolo.
  \item[compute] Il nodo di calcolo con GPU. Ha il container ma non l'ambiente
        installato sull'host.
  \item[apptainer] Il container che fornisce l'ambiente di esecuzione.
  \item[setup] Lo script che, all'interno del container, prepara l'ambiente.
\end{description}

Il fatto che il nodo di login non abbia il runtime è una fonte ricorrente di
confusione, perché il modo naturale di verificare qualcosa (``eseguo un
\texttt{python -c} per controllare'') non funziona lì.

\subsection{Il container: \texttt{exec} e non \texttt{run}}

Il container viene invocato con \texttt{apptainer exec}, non con
\texttt{apptainer run}. La ragione è concreta: \texttt{run} passa attraverso lo
\emph{runscript} del container, che manipola gli argomenti, e questa
manipolazione corrompe la lista degli argomenti quando si invoca
\texttt{python -c} con codice che contiene virgolette e spazi. Il sintomo è un
errore di sintassi Python su codice sintatticamente valido, che è una
diagnostica particolarmente inutile.

L'invocazione è centralizzata nella funzione \texttt{run\_py} di
\texttt{cluster/\_lib.sh}, che costruisce il comando con dieci direttive
\texttt{--env} per propagare le variabili d'ambiente necessarie. La
centralizzazione è deliberata: se ogni script costruisse la propria
invocazione, l'insieme delle variabili propagate divergerebbe fra script e i
guasti sarebbero specifici del punto di ingresso.

\section{Il regime offline}

I nodi di calcolo non hanno accesso alla rete. Questo vincolo è più stringente
di quanto sembri, perché le librerie moderne di machine learning tentano
connessioni di rete per motivi non ovvi: verificare la presenza di aggiornamenti
del modello, risolvere un identificativo di dataset, inizializzare la
telemetria di un servizio di logging.

L'unico percorso in cui la rete è disponibile è \texttt{setup.sh}, che scarica
in anticipo modello e dati nella cache locale. Tutti i job successivi girano
offline.

La configurazione dell'ambiente offline è raccolta nella funzione
\texttt{export\_offline\_env}, che esporta:

\begin{center}
\begin{tabular}{ll}
\toprule
Variabile & Effetto \\
\midrule
\texttt{HF\_HUB\_OFFLINE} & disattiva le chiamate all'hub dei modelli \\
\texttt{TRANSFORMERS\_OFFLINE} & idem, a livello di libreria \\
\texttt{HF\_DATASETS\_OFFLINE} & idem, per i dataset \\
\texttt{WANDB\_MODE=offline} & registrazione locale, nessun invio \\
\texttt{WANDB\_DISABLE\_WEAVE} & disattiva un componente accessorio \\
\texttt{WANDB\_SILENT} & riduce l'output \\
\texttt{PYTHONUNBUFFERED} & log immediati, essenziale per diagnosticare i job \\
\texttt{HF\_HOME} & radice della cache (\texttt{\$HOME/.cache/huggingface}) \\
\texttt{HF\_HUB\_CACHE} & cache specifica dell'hub \\
\bottomrule
\end{tabular}
\end{center}

\textbf{Regola critica}: \texttt{export\_offline\_env} va invocata
\emph{prima di qualunque processo Python}. Le librerie leggono queste variabili
al momento dell'import, non all'uso: esportarle dopo il primo import non ha
effetto. La sezione~\ref{sec:bug-offline-order} documenta il caso in cui questa
regola era violata.

\section{La catena di job}

La quota assegnata consente \textbf{una sola allocazione} alla volta. Non è
possibile lanciare più esperimenti in parallelo; le celle sperimentali devono
essere eseguite in sequenza.

Questo ha richiesto un meccanismo di concatenazione: uno script
(\texttt{job\_chain}) mantiene una coda persistente di celle da eseguire, e a
ogni conclusione di job accoda il successivo. Il meccanismo è basato su
\emph{tick} periodici, usa \texttt{flock} per garantire l'accesso esclusivo alla
coda (due tick concorrenti corromperebbero lo stato) e implementa un ritentativo
sulle interrogazioni a \texttt{sacct}, che sotto carico può non rispondere.

Il monitoraggio è realizzato da \texttt{python3 -m src.utils.chain\_monitor}, con
una modalità \texttt{monitor --all} che mostra lo stato di tutta la coda.

\section{Sette guasti reali}

Ciò che segue è la parte utile del capitolo.

\subsection{Guasto 1: percorso dei dati divergente fra test e produzione}

Lo script di preparazione dei dati verificava l'esistenza del percorso
\texttt{data/aslg\_pc12\_train}. La produzione, però, usa
\texttt{data/aslg\_pc12}. Il percorso testato non esisteva.

La conseguenza: il controllo di presenza della cache falliva sempre, quindi lo
script tentava il \emph{download} del dataset. Su un nodo offline, il download
fallisce. Il sintomo era dunque un errore di rete durante una fase che avrebbe
dovuto essere puramente locale, con la diagnostica che indirizzava verso la
configurazione offline invece che verso un percorso sbagliato.

Lezione: un controllo di esistenza su un percorso che non è quello usato in
produzione è peggio di nessun controllo, perché produce un ramo di fallback che
non doveva essere raggiunto.

\subsection{Guasto 2: ambiente offline esportato dopo il primo Python}
\label{sec:bug-offline-order}

Nello script di valutazione, \texttt{export\_offline\_env} era invocata alla riga
$177$. Il primo processo Python veniva però lanciato alla riga $86$, per
risolvere la directory di uscita.

Quel processo girava dunque \emph{senza} le variabili offline, tentava una
connessione, e in assenza di rete si bloccava fino al timeout. Il job appariva
sospeso in una fase di inizializzazione banale.

Lezione: l'ordine delle esportazioni d'ambiente rispetto al primo import è un
vincolo di correttezza, non di stile. La contromisura strutturale è invocare la
configurazione d'ambiente come primissima istruzione dello script, prima di
qualunque logica.

\subsection{Guasto 3: cambio di firma in Transformers 5.3.0}
\label{sec:bug-weave}

Questo è il guasto tecnicamente più interessante.

Transformers $5{,}3{,}0$ ha modificato una funzione interna di rilevamento dei
pacchetti, \texttt{\_is\_package\_available}, che ora restituisce una
\emph{tupla} invece di un booleano. Il codice chiamante valutava il risultato in
un contesto booleano. Ma in Python
\begin{equation}
  \texttt{bool((False, None))} \;=\; \texttt{True}
\end{equation}
perché una tupla non vuota è sempre veritiera, indipendentemente dal suo
contenuto.

Il risultato: il codice concludeva che il pacchetto di telemetria \emph{era}
disponibile, tentava di importarlo, e l'import falliva. Poiché la catena di
import passava per \texttt{trl.trainer.grpo\_trainer}, l'effetto era che
\textbf{l'intero modulo del trainer GRPO diventava non importabile}, con sette
test in rosso e un errore che non menzionava né TRL né il GRPO.

Lezione doppia. Primo: verificare esplicitamente il tipo dei valori restituiti
dalle API interne di librerie di terze parti, che non seguono garanzie di
compatibilità. Secondo, e più generale: la veridicità implicita in Python
trasforma un cambio di tipo in un cambio di comportamento silenzioso.
L'esportazione di \texttt{WANDB\_DISABLE\_WEAVE} nella lista della sezione
precedente è la contromisura applicata.

\subsection{Guasto 4: la coda invertita}

La funzione di ricostruzione della coda dei job effettuava un \emph{prepend},
cioè inseriva in testa, invece di un \emph{append}.

L'effetto: ricostruendo una coda che doveva essere
$[\texttt{train}, \texttt{eval}]$ si otteneva $[\texttt{eval}, \texttt{train}]$.
La valutazione girava prima dell'addestramento, su un checkpoint inesistente o
obsoleto.

Il guasto è insidioso perché entrambi i job \emph{riuscivano}: la valutazione
trovava un checkpoint (quello precedente) e produceva metriche plausibili.
L'errore si manifestava come risultati sperimentali che non corrispondevano alla
configurazione dichiarata, che è la peggiore classe di guasto in un contesto
sperimentale.

Lezione: nelle strutture dati che rappresentano un ordine di esecuzione,
l'orientamento dell'inserimento va verificato con un test esplicito. Un
commento non basta.

\subsection{Guasto 5: celle di sola valutazione lanciate come addestramento}
\label{sec:bug-evalonly}

Le celle \texttt{baseline/*} sono celle di \emph{sola valutazione}: misurano il
modello base, senza addestrarlo. Le loro configurazioni non contengono quindi la
sezione relativa all'addestramento, e in particolare non contengono la chiave
\texttt{output\_dir}.

Lanciandone una con lo script di addestramento, il risultato era:
\begin{center}
\texttt{KeyError: 'output\_dir'}
\end{center}

L'errore in sé è corretto: la configurazione non ha quella chiave, perché non
deve averla. Il problema era \emph{quando} si manifestava. La validazione
avveniva dopo il caricamento di Unsloth e del modello, quindi il job occupava una
GPU per minuti prima di fallire con un errore di configurazione diagnosticabile
in millisecondi. In un regime a una sola allocazione, ogni minuto sprecato è un
esperimento non eseguito.

La correzione è su due livelli, per difesa in profondità:

\begin{enumerate}
  \item nel punto di ingresso \texttt{src/training/\_\_main\_\_.py}, una guardia
        posta \textbf{prima} dell'import di Unsloth. Riconosce le celle di sola
        valutazione, fallisce in circa un secondo, e il messaggio d'errore indica
        lo script corretto da usare (\texttt{cluster/eval.sh}). L'uscita avviene
        con codice $2$, distinguibile da un fallimento generico;
  \item nella funzione principale di \texttt{src/training/grpo\_t2g\_train.py},
        la stessa guardia ripetuta, con messaggi distinti per i due casi: cella
        di sola valutazione contro configurazione genuinamente incompleta.
\end{enumerate}

La copertura è garantita da dieci test in
\texttt{tests/test\_train\_entrypoint\_guard.py}, che verificano sia
l'invariante diretta (una cella di sola valutazione viene rifiutata) sia
l'invariante inversa (una cella addestrabile \emph{non} viene rifiutata). Il
secondo gruppo è quello che conta di più: una guardia troppo aggressiva
rifiuterebbe esperimenti validi, e sarebbe un guasto peggiore di quello che
corregge.

Lezione, che è la stessa del guasto 2 in forma generale: \textbf{validare prima
di allocare}. Ogni controllo che può essere eseguito senza caricare il modello
deve essere eseguito prima di caricarlo. È anche il principio codificato nel
contratto degli obiettivi ausiliari
(sezione~\ref{sec:contratto-ausiliari}), dove
\texttt{resolve\_auxiliary\_config} valida prima del caricamento.

\subsection{Guasto 6: la fusione dell'adattatore SFT mai salvata su disco}
\label{sec:bug-sft-merge-eval}

Il più recente, e il più rilevante per l'interpretazione dei risultati del
Capitolo~\ref{cap:risultati}. Riguarda ogni cella \texttt{sft-grpo/*}: quelle
in cui una fase di SFT precede il GRPO nella stessa run.

\texttt{src/models/model\_loader.py} fonde l'adattatore SFT nel modello base
\textbf{solo in memoria}, prima di agganciare un nuovo LoRA per la fase GRPO.
La fusione non viene mai scritta su disco: \texttt{trainer.save\_model()},
applicato al modello PEFT risultante, salva soltanto il \emph{delta}
dell'adattatore GRPO, calcolato rispetto a una base+SFT che non viene
registrata da nessuna parte come oggetto a sé.

\texttt{src/training/eval\_t2g.py::load\_model\_for\_eval} ricostruiva il
modello come \texttt{base grezzo + adattatore GRPO}, senza sapere che una fase
SFT fosse mai avvenuta. Il contributo dell'SFT --- la fase che consuma la
maggior parte del tempo di calcolo della run --- spariva silenziosamente da
ogni valutazione di ogni cella \texttt{sft-grpo}, senza errore, senza validità
alterata: un numero plausibile, calcolato sul modello sbagliato.

\subsubsection{Perché non si notava}

L'adattatore SFT \emph{sopravvive} sul disco: \texttt{grpo\_t2g\_train.py}
addestra (o copia) la fase SFT in
\texttt{<run\_dir>/sft\_pretrain/final/}, un fratello della cartella
\texttt{final/} del GRPO, proprio per rendere la run autocontenuta. Il difetto
non era un dato mancante, ma un percorso mai consultato: il codice di
valutazione non sapeva che quella cartella andasse cercata.

È stato trovato rileggendo il codice del caricamento del modello e
confrontandolo con il layout reale delle run sul cluster (\texttt{find} sulla
struttura di una cella \texttt{sft-grpo} completata), non da un sintomo
osservabile nei log o nelle metriche.

\subsubsection{La correzione}

Una funzione di individuazione,
\texttt{\_sft\_pretrain\_adapter\_sibling(checkpoint\_path)}, cerca il fratello
\texttt{sft\_pretrain/final/} della cartella del checkpoint GRPO (per run
complete e per checkpoint intermedi \texttt{checkpoint-N}, non solo per
\texttt{final}). Quando esiste, \texttt{load\_model\_for\_eval} lo fonde nel
modello base \textbf{prima} di applicare l'adattatore GRPO, ripristinando
l'ordine di composizione realmente usato in addestramento. Sei test in
\texttt{tests/test\_sft\_grpo\_eval\_composition.py} coprono il rilevamento:
il caso positivo, l'assenza di fase SFT (celle \texttt{grpo/*} pure), una fase
SFT incompleta (mai trattata come un merge avvenuto), e un checkpoint
intermedio.

\subsubsection{L'esito: la prima ipotesi}
\label{sec:sft-merge-non-rimisurato}

Il paragrafo precedente, nella stesura originale di questa relazione,
lasciava aperte due ipotesi e rimandava lo scioglimento a una
ri-valutazione a costo pressoché nullo della cella \texttt{sft-grpo} già
addestrata. Nel frattempo lo storage del cluster è stato azzerato per
liberare spazio, quindi l'esperimento eseguito non è stato la
ri-valutazione dello stesso checkpoint, ma una campagna fresca (\texttt{max\_steps}
$2\,000 \to 5\,000$, valutazione doppia) con la correzione già in codice fin
dall'addestramento. La domanda a cui risponde è comunque la stessa: un
checkpoint \texttt{sft-grpo} composto correttamente si avvicina al soffitto
dell'SFT, o resta al livello del vecchio $0{,}5114$?

\textbf{Si conferma la prima ipotesi.} La cella \texttt{sft-grpo/few-shot}
rivalutata ottiene ROUGE-L $0{,}9628$ ed exact match $0{,}7296$ (few-shot
nativo) — contro $0{,}9681$/$0{,}7662$ dell'SFT da sola nella stessa
campagna (zero-shot nativo): un divario residuo di circa $0{,}01$ di
ROUGE-L, non di $0{,}46$. Il numero storico $0{,}5114$ misurava per gran parte
\texttt{base grezzo + adattatore GRPO}, cioè esattamente il guasto appena
descritto, non l'effetto del reinforcement learning sulla policy SFT.

Questo non annulla il meccanismo di collasso del supporto
(sezione~\ref{sec:supporto}): l'$84{,}8\%$ di gruppi a varianza nulla è una
proprietà misurata direttamente sui rollout, indipendente da come viene
poi ricomposto il checkpoint per l'eval, e resta un fatto reale. Quello che
cambia è la sua portata: un meccanismo che azzera il gradiente su gran
parte dei gruppi produce, una volta il checkpoint valutato correttamente,
una deriva piccola (circa un punto di ROUGE-L) e non il crollo di quaranta
punti riportato dall'Osservazione~5 originale. Il meccanismo era la
spiegazione corretta di \emph{una parte contenuta} del divario; il resto
era il guasto stesso. Il dettaglio con la tabella di confronto è
riportato nel Capitolo~\ref{cap:risultati}, sezione~\ref{sec:osservazione5}.

\subsection{Guasto 7: la catena della pipeline degli obiettivi ausiliari}
\label{sec:bug-ausiliari}

Non un guasto ma una catena: nove difetti incontrati uno dopo l'altro mettendo
in esercizio gli obiettivi del Capitolo~\ref{cap:ausiliari}, ognuno visibile
solo dopo aver corretto il precedente. Sono raccolti insieme perché hanno una
causa di fondo comune, discussa sotto.

\begin{center}
\small
\begin{tabular}{p{0.3cm}p{5.2cm}p{5.6cm}l}
\toprule
\# & Sintomo & Causa & Tipo \\
\midrule
1 & crash del collatore al primo batch & classe base \texttt{DataCollatorForLanguageModeling} di Transformers (MLM) invece dell'omonima di TRL & rumoroso \\
2 & \texttt{NotImplementedError} sui logits & Unsloth restituisce un segnaposto \texttt{EmptyLogits}; il termine strutturato lo toccava senza averne bisogno & rumoroso \\
3 & termine strutturato mai applicato, $0$ righe valutate per $800$ passi & la preparazione del dataset compilata da Unsloth rimuove tutte le colonne originali, \texttt{gold\_gloss} compresa & silenzioso \\
4 & crash nella mappatura dei bersagli & il grafo costruito non veniva mai passato al trainer & rumoroso \\
5 & tensori su dispositivi diversi & i buffer della loss strutturata restavano su CPU & rumoroso \\
6 & logits vuoti nonostante \texttt{UNSLOTH\_RETURN\_LOGITS=1} & la variabile torna a \texttt{0} durante \texttt{Trainer.train()} & rumoroso \\
7 & massa ammessa mai applicata, $0$ posizioni per $100$ passi & controllo di fine sequenza sull'ultima posizione, occupata dall'a capo dopo \texttt{<|im\_end|>} & silenzioso \\
8 & loss di validazione $+\infty$, arresto a $800$ passi & una riga con percorso gold non rappresentabile rendeva infinita la media del batch & silenzioso \\
9 & controllo mescolato che vince sul grafo vero & la permutazione delle destinazioni spegneva il termine nel $98\%$ dei passi & silenzioso \\
\bottomrule
\end{tabular}
\end{center}

\subsubsection{La causa di fondo: il codice sul disco non è quello che gira}

Tre dei nove difetti (2, 3 e 6), i più costosi da diagnosticare, nascono dallo
stesso fraintendimento: aver letto il sorgente installato delle librerie e
averlo creduto il codice eseguito. Unsloth non si limita a modificare il modello: al
momento dell'import compila e sostituisce metodi interi del trainer di TRL. La
preparazione del dataset effettivamente eseguita, ad esempio, è materializzata
in \texttt{unsloth\_compiled\_cache/UnslothSFTTrainer.py} ed è diversa da quella
del pacchetto TRL; lo stesso trainer compilato non eredita dall'\texttt{SFTTrainer}
di TRL ma da una sua classe base comune, per cui il \texttt{compute\_loss} letto
nel sorgente di TRL non viene mai eseguito. Il difetto 6 è la forma più pura del
problema: la variabile d'ambiente veniva impostata correttamente, prima
dell'import, e verificata nei log; il suo valore al momento del passaggio in
avanti, letto con una sonda nel ciclo di addestramento vero, era \texttt{0}.
Chiamare lo stesso \texttt{compute\_loss} fuori da \texttt{Trainer.train()}
restituiva invece logits reali.

La diagnosi è arrivata solo quando si è smesso di ragionare sul sorgente e si è
interrogato il codice vivo sul cluster: \texttt{inspect.getsource} sui metodi
effettivamente legati all'istanza del trainer e una sonda sul valore delle
variabili al momento della chiamata. La correzione del difetto 6 non dipende
dall'aver individuato il meccanismo interno che azzera la variabile: una
callback la riafferma all'inizio di ogni passo, e poiché il passaggio in avanti
compilato la legge a ogni chiamata, vince comunque.

\subsubsection{Perché i test non li avevano visti}

I test degli obiettivi ausiliari chiamavano \texttt{compute\_loss} con un
dizionario di input costruito a mano, che conteneva già \texttt{gold\_gloss},
bersagli con fine sequenza nell'ultima posizione, tensori già sul dispositivo
giusto e percorsi sempre rappresentabili. Ciascuna di queste comodità
nascondeva uno dei difetti. Ogni correzione è accompagnata da un test di
regressione che riproduce la condizione reale: il segnaposto dei logits che
solleva un'eccezione a ogni accesso, l'input senza \texttt{gold\_gloss}, il
token finale dopo la fine sequenza, la riga con percorso non rappresentabile, il
controllo che deve lasciare valutabili i percorsi di training.

\subsubsection{Gli indicatori che hanno reso visibili i silenziosi}

I quattro difetti silenziosi sono emersi solo perché il trainer registra, a ogni
passo, quante righe o posizioni il termine ausiliario ha effettivamente
valutato, e (per il termine strutturato) un codice che dice \emph{perché} un
passo non ne ha valutata nessuna: rampa non iniziata, bersaglio assente, stati
nascosti assenti, righe ineleggibili, percorsi non rappresentabili. Un
contatore fermo a zero per centinaia di passi è stato ogni volta il primo
indizio; il codice di motivazione ha trasformato ogni ripartenza in una misura
invece che in un tentativo.

\section{Il filo comune}

I sette guasti hanno cause diverse ma condividono una struttura: nessuno era un
errore di matematica, tutti riguardavano l'\emph{ordine} delle operazioni o
l'\emph{ambiente} in cui avvenivano.

\begin{center}
\begin{tabular}{ll}
\toprule
Guasto & Categoria \\
\midrule
1. percorso dati divergente & ambiente: percorso testato $\neq$ usato \\
2. offline dopo il primo Python & ordine: esportazione dopo l'import \\
3. tupla veritiera & ambiente: API interna di terze parti \\
4. coda invertita & ordine: prepend invece di append \\
5. validazione dopo l'allocazione & ordine: controllo dopo il costo \\
6. fusione SFT mai salvata & ambiente: composizione del checkpoint non registrata \\
7. pipeline degli obiettivi ausiliari & ambiente: codice di terze parti sostituito a runtime \\
\bottomrule
\end{tabular}
\end{center}

Quattro su sei producevano fallimenti \emph{plausibili}: il job girava, non
segnalava nulla di anomalo, e restituiva numeri. Su un progetto sperimentale
questa è la categoria di guasto più costosa, perché contamina i risultati senza
lasciare traccia nei log. Il guasto 6 appartiene alla stessa categoria in modo
particolarmente puro: non un job che falliva in modo silenzioso, ma un job che
\emph{riusciva} nel senso pieno del termine, producendo un numero corretto per
un modello diverso da quello dichiarato. Il guasto 7 ne aggiunge quattro della
stessa specie: un termine ausiliario che non si applicava, un altro che non si
applicava per una ragione diversa, un arresto anticipato con numeri plausibili e
un controllo che controllava un'altra cosa.

La contromisura adottata in tutti i casi è la stessa e vale come conclusione del
capitolo: \textbf{fallire presto e in modo rumoroso}, e codificare ogni
invariante scoperta in un test che ne impedisca il ritorno.
